\chapter{什么是 AI、机器学习与科学推断}

\section{学习目标}

学完本章后，你应该能够：

\begin{itemize}
    \item 区分规则方法、统计模型和机器学习模型；
    \item 理解预测任务与科学解释任务并不完全相同；
    \item 认识为什么机器学习必须和 baseline 比较，而不是凭“看起来高级”就直接采用；
    \item 用一个极小的恒星分类例子建立最小分类流程。
\end{itemize}

\section{AI、机器学习与科学问题}

在今天的语境里，AI 是一个非常宽的词。对本教材来说，我们更关心的是其中与科研最直接相关的一部分：机器学习。

\begin{itemize}
    \item \textbf{规则方法}：由人手工写出明确规则；
    \item \textbf{统计模型}：利用参数化模型描述数据关系；
    \item \textbf{机器学习}：让模型通过数据学习决策边界或映射关系。
\end{itemize}

这三类方法并不是互相排斥的。在科学工作中，它们常常应该协同使用，而不是被当作彼此替代的“阵营”。

\section{预测与解释不是同一件事}

机器学习很擅长做预测，但科学研究除了预测，还关心解释：

\begin{itemize}
    \item 为什么这些天体会分成不同类别？
    \item 哪些特征真正携带了物理信息？
    \item 一个模型分数更高，是否就意味着它更接近物理机制？
\end{itemize}

因此，本教材会不断提醒一件事：模型的任务是帮助你组织和利用数据，而不是直接替代物理解释。

\section{为什么需要 baseline}

在配套 notebook 中，我们先不用任何机器学习库，而是写一个基于颜色和绝对星等的规则分类器：

\begin{examplebox}[规则 baseline]
\begin{lstlisting}[style=pythonstyle]
def classify_star_rule(bp_rp, absolute_g_mag):
    if absolute_g_mag > 10.0:
        return "white_dwarf"
    if absolute_g_mag < 2.0 and bp_rp > 1.2:
        return "giant"
    return "main_sequence"
\end{lstlisting}
\end{examplebox}

这一步很重要，因为它让我们先有一个可以解释、可以比较的出发点。后面如果一个复杂模型只比这个 baseline 好一点点，就说明问题可能并不需要特别复杂的方法。

\section{一个最小恒星分类示例}

配套 notebook 使用一个小型 Gaia 风格教学样本，把每颗恒星映射到：

\begin{itemize}
    \item 颜色指数；
    \item 绝对星等；
    \item 恒星类别标签。
\end{itemize}

这其实已经构成了一个最小的监督学习任务：

\begin{itemize}
    \item 特征：颜色指数、绝对星等；
    \item 标签：主序星、巨星、白矮星；
    \item 目标：根据特征预测标签。
\end{itemize}

\section{AI 助手如何帮助你}

在这一章，AI 助手最适合做的是：

\begin{itemize}
    \item 帮你解释一个 baseline 为什么合理或不合理；
    \item 把规则分类器整理成更清楚的代码；
    \item 对比不同方法的优缺点。
\end{itemize}

但你需要自己判断：某个模型为什么有效、它利用了什么特征、它是否真的服务于科学问题。

\section{本章小结}

机器学习不是“把数据丢进黑箱”，也不是“比传统方法更时髦所以一定更好”。它是一组建立在 baseline、特征、训练数据和评估标准之上的方法。真正重要的是：这些方法是否帮助你更清楚地组织证据、提出判断并理解局限。
